\providecommand{\fqeCalibFigDir}{FQE_calibration_neurips/paper_import_bundle/figures}

\section{Experiments}
\label{sec:experiments}

We evaluate Bellman calibration as a post-hoc value-space wrapper for fitted
off-policy value predictors. Each baseline first estimates \(Q(s,a)\), inducing
\(\widehat V(s)=\sum_a\pi(a\mid s)\widehat Q(s,a)\). Calibration then learns a
scalar map \(g\) and evaluates \(g(\widehat V)\). The calibration map is fit by
importance-weighted fitted value iteration,
\[
  g_{t+1}\in\arg\min_g\sum_i \rho_i
  \left[g(\widehat V_i)-\{R_i+\gamma g_t(\widehat V'_i)\}\right]^2,
  \qquad
  \rho_i=\pi(A_i\mid S_i)/\mu(A_i\mid S_i),
\]
where \((\widehat V_i,\widehat V'_i)\) are out-of-fold predictions for
cross-calibrated rows. Temporal reward-shift rows instead train the raw value
model on old-regime behavior data and fit \(g\) on a small held-out batch from
the recent current regime. The same-small-current retrain comparator uses the
same learner with a prespecified prediction clip to avoid numerical explosions,
and is reported as a comparator rather than promoted as a calibrated method. All
promoted claims are evaluated on independent current or same-regime test data
and independent oracle rollouts.

The final audit uses 100 independent replications for all core experiment
rows. We report ratios to the matched uncalibrated baseline for true value
function MSE and binned Bellman calibration error; values below one are better.
Brier-style Bellman outcome errors are saved as diagnostics but are not used for
promotion.

\begin{figure*}[t]
  \centering
  \includegraphics[width=0.82\linewidth]{\fqeCalibFigDir/calibration_story_compact.pdf}
  \caption{\textbf{Bellman calibration improves value accuracy and calibration
  in score-correctable regimes.} Each bar is a ratio to the matched
  uncalibrated baseline. The organic rows show affine/misspecified linear FQE,
  finite-iteration random-feature FQE, and temporal reward shift with recent
  calibration. Mechanism rows illustrate that linear calibration is enough for
  affine distortion while isotonic is better for monotone saturation.}
  \label{fig:calibration-story}
\end{figure*}

\begin{table*}[t]
\centering
\caption{\textbf{Representative audited Bellman calibration results.}
Ratios are relative to the matched uncalibrated baseline. Win rates are
seed-matched true-\(V\) MSE win rates unless otherwise noted.}
\label{tab:evidence-audit-main}
\scriptsize
\setlength{\tabcolsep}{4pt}
\begin{tabular}{@{}p{0.43\linewidth}lrrr@{}}
\toprule
Regime and learner & Calibrator & Rel. \(V\) MSE & Rel. cal. err. & Win rate \\
\midrule
Affine misspec., linear FQE & Linear & 0.750 & 0.631 & 0.94 \\
Affine misspec., linear FQE & Isotonic & 0.734 & 0.764 & 0.94 \\
Finite iteration, RF FQE \(K=2\) & Linear & 0.781 & 0.664 & 0.95 \\
Finite iteration, RF FQE \(K=4\) & Isotonic & 0.776 & 0.736 & 0.94 \\
Temporal reward shift, RF FQE & Linear recent & 0.630 & 0.717 & 0.97 \\
Temporal reward shift, RF FQE & Isotonic recent & 0.627 & 0.916 & 0.98 \\
\addlinespace[2pt]
\multicolumn{5}{@{}l}{\emph{Mechanism diagnostics, not organic learner claims}}\\
Affine distortion mechanism & Linear & 0.430 & 0.169 & 1.00 \\
Monotone saturation mechanism & Isotonic & 0.375 & 0.022 & 1.00 \\
\bottomrule
\end{tabular}
\end{table*}

Figure~\ref{fig:calibration-story} and Table~\ref{tab:evidence-audit-main}
show aligned improvements in value accuracy and Bellman calibration. In the
affine/misspecified linear-FQE regime, linear calibration is sufficient:
relative true-\(V\) MSE falls to about \(0.75\) and Bellman calibration error to
about \(0.63\). In the finite-iteration regime, the raw random-feature FQE
scores remain strongly rank-informative but are mis-scaled; calibration reduces
true-\(V\) MSE to about \(0.78\) and Bellman calibration error to
\(0.66\)--\(0.74\). In the temporal reward-shift regime, a model trained on
old-reward data is recalibrated with a small recent current-regime batch,
reducing current-regime true-\(V\) MSE by about \(37\%\) and Bellman calibration
error by about \(28\%\) for linear calibration.

As a benchmark check, we take the official Deep OPE FQE-L2 scores for
\texttt{cheetah\_run} from DeepMind's RL Unplugged suite as a fixed fitted-value
baseline, and only fit the post-hoc one-dimensional maps. Across three official
learning curves, the raw FQE-L2 scores are rank-informative but mis-scaled
(Spearman \(0.67\)--\(0.88\)). We calibrate on one subset of policy snapshots
and evaluate absolute OPE error on disjoint snapshots, using three fixed splits
within each curve. Linear and isotonic maps reduce error on every split; averaged
over curves, relative OPE error is \(0.39/0.36\) for policies
\(0\)--\(3\) calibrated and \(4\)--\(7\) evaluated, \(0.20/0.19\) for
even/odd splits, and \(0.36/0.31\) for policies \(0\)--\(4\) calibrated and
\(5\)--\(7\) evaluated (linear/isotonic;
Figure~\ref{fig:cheetah-policy-value}). No base value learner is refit.

\begin{figure}[t]
  \centering
  \includegraphics[width=0.68\linewidth]{\fqeCalibFigDir/rlu_cheetah_policy_calibration_compact.pdf}
  \caption{\textbf{Deep OPE/RL Unplugged \texttt{cheetah\_run}.}
  Absolute OPE error after post-hoc value calibration of official Deep OPE
  FQE-L2 scores, relative to the raw score. Bars average over three official
  learning curves; whiskers show the range. The dashed line is the raw baseline,
  and row labels give calibration and held-out evaluation policy indices.}
  \label{fig:cheetah-policy-value}
\end{figure}
